\section{Data and computations}\label{app:data}

\subsection{The complex \texorpdfstring{$N_5$}{N\_5}}\label{ss:N5}
Encoding $L_{5,\mathrm{PL}}$ in the chart of \S\ref{ss:explicit} and performing the unit
cancellations gives the reduced twisted complex $N_5$ on the generators $0,1,\dots,22$.  Its
$\iota_\circ$ generators are $3,4,9,10,13,14,19,20$, the others carry $\iota_\bullet$, and
\begin{equation}\label{eq:N5}
\begin{aligned}
\delta_{N_5}={}&1\xrightarrow{M^2}0+2\xrightarrow R1+3\xrightarrow M2+4\xrightarrow L3
+4\xrightarrow M5+5\xrightarrow R6+6\xrightarrow{M^2}7+8\xrightarrow R7\\
&+9\xrightarrow M8+10\xrightarrow L9+10\xrightarrow M11+12\xrightarrow R11
+13\xrightarrow M12+14\xrightarrow L13+14\xrightarrow M15\\
&+16\xrightarrow R15+17\xrightarrow{M^2}16+18\xrightarrow R17+19\xrightarrow M18
+20\xrightarrow L19+20\xrightarrow M21+21\xrightarrow R22.
\end{aligned}
\end{equation}
No arrow ends where an arrow with the same face label begins, so $\delta_{N_5}^2=0$.  The complex is a
single arc with ends $0$, at $(0,0)$, and $22$, at $(\pi,0)$.  The encoding of a numerical trace of
Smith's perturbed curve at $q=5$, computed as in Remark~\ref{rem:analytic}, is strictly isomorphic to
$N_5$.  Its pairing dimensions with $E_{-1/2}$ and $E_{-3/4}$ are $5$ and $15$; the first agrees with Smith's
rank $5$ for this decomposition \cite[proof of Thm.~5.9]{Smith}.

The smoothing used in Theorem~\ref{thm:structure-full} is the four-arrow element
\begin{equation}\label{eq:b5}
 b_5=8\xrightarrow R7+16\xrightarrow R15+8\xrightarrow R15+16\xrightarrow R7 ,
\end{equation}
which exchanges the targets of the arrows $8\xrightarrow R7$ and $16\xrightarrow R15$.  The complex
$(N_5,b_5)$ consists of the arc $0,1,\dots,7,16,17,\dots,22$ and the closed component
$8,9,\dots,15$, whose word $9\xrightarrow M8$, $10\xrightarrow L9$, $10\xrightarrow M11$,
$12\xrightarrow R11$, $13\xrightarrow M12$, $14\xrightarrow L13$, $14\xrightarrow M15$,
$8\xrightarrow R15$ is the doubled word of $c$.  Adding $\kappa_5=12\xrightarrow R11+0\xrightarrow
R11$ opens the closed component between $11$ and $12$ and attaches $11$ to the end $0$; the result
is the single arc $22,21,\dots,16,7,6,\dots,0,11,10,9,8,15,14,13,12$.  Its pairing dimensions with
$E_{-1/2}$ and $E_{-3/4}$ are $7$ and $21$, the ranks of reduced Khovanov homology of $P(-2,3,5)$ and
$K_{-3/4}(5)$.

\subsection{Corner terms}\label{ss:corner-terms}
The search of Remark~\ref{rem:corner-term} finds the following one-arrow replacements carrying
$(N_q,b)$ to a complex strictly isomorphic to $\mathrm{BN}_q$, written as elements of
$\operatorname{End}(N_q)$ (the first term is the deleted arrow):
\begin{center}
\begin{tabular}{lll}
\toprule
object & first solution & second solution\\
\midrule
$(N_7,b_{S_{18}})$ & $\kappa_7=0\xrightarrow R19+36\xrightarrow R19$ & $22\xrightarrow R3+36\xrightarrow R3$\\
$(N_7,b_{S_{25}})$ & $\kappa_7=0\xrightarrow R19+36\xrightarrow R19$ & $4\xrightarrow R3+36\xrightarrow R3$\\
$(N_5,b_5)$ & $\kappa_5=12\xrightarrow R11+0\xrightarrow R11$ & $8\xrightarrow R15+0\xrightarrow R15$\\
\bottomrule
\end{tabular}
\end{center}
The arrows $22\xrightarrow R3$ and $8\xrightarrow R15$ are arrows of $b_{S_{18}}$ and $b_5$ rather
than of $\delta_{N_q}$, so the second solutions in the first and third rows delete an arrow of the
smoothing.  The term $\kappa_7$ is the only one that serves both smoothings at $q=7$.  The
analogous search with $\widetilde{\mathrm{BN}}(T_q)$ in place of $\mathrm{BN}_q$, which is the
comparison for $t>0$ (Remark~\ref{rem:sign-positive}), finds for both smoothings at $q=7$ the terms
$2\xrightarrow M3+18\xrightarrow{M^2}3$ and $20\xrightarrow M19+18\xrightarrow{M^2}19$, and at $q=5$
the terms $10\xrightarrow M11+22\xrightarrow{M^2}11$ and $14\xrightarrow M15+22\xrightarrow{M^2}15$.
The exact edit distance from $N_q$ to $\mathrm{BN}_q$ is $2$ for $q=5,7$, with two minimal edits in
each case; at $q=7$ both consist of arrows labelled $R$ and touch the end $36$.

\subsection{Finite polygon counts and candidate supports}\label{sec:compute}
For $q=5,7,11,13$ we record the output of a finite combinatorial test on the piecewise-linear
curves $R_{\mathrm{PL}}=\Rnat(\Qh{-1/2})$ and $L_{q,\mathrm{PL}}$; at $q=7$ it is superseded by the exact
calculation of Proposition~\ref{prop:q7-typeD}.  De~Silva, Robbin and Salamon
identify holomorphic strips with combinatorial lunes for transverse embedded loops on a surface
\cite{dSRS}; their hypotheses do not cover immersed curves in the orbifold pillowcase, nor the higher
polygons of a bounding-cochain deformation.  The statements labelled Computation below record outputs of
this test as data.

\subsubsection{The finite winding test}\label{ss:polygons}
For intersection points $x,y$ choose an arc $\alpha$ on $R_{\mathrm{PL}}$ from $x$ to $y$ and an arc
$\beta$ on $L_{q,\mathrm{PL}}$ from $y$ to $x$, and lift $\alpha\cdot\bar\beta$ to the torus.  The pair is a
\emph{candidate bigon} when (a) the lifted loop is null-homologous, (b) its winding function vanishes at
the four corners, (c) that function is nonnegative on the sampled complementary regions, and (d) the
local sectors at $x$ and $y$ are convex.  The same conditions at every vertex define candidate triangles
and quadrilaterals; we call their conjunction the \emph{low-order test}.  The matrices use one
fixed lift of $R_{\mathrm{PL}}$ and sum over the two torus preimages of a self-intersection of
$L_{q,\mathrm{PL}}$.  The arcs $\alpha$ and $\beta$ range over explicitly listed subarcs of the two curves,
each traversed at most once, and no geometric bound shows that these subarcs exhaust all candidates.  We
write $A_q$ for the set of self-intersections of $L_{q,\mathrm{PL}}$ that are vertices of a candidate
triangle or quadrilateral found in this way.  The parameters $(\varepsilon_L,\varepsilon_R,r_R)$, namely
the connector amplitude of $L_{q,\mathrm{PL}}$ and the perturbation amplitude and pinch parameter of
$R_{\mathrm{PL}}$, are $(0.05,0.10,0.25)$ at $q=5$ and $(0.05,0.16,0.40)$ at $q=7,11,13$
unless stated otherwise.  With $C_q$ the $\F_2$-vector space on the transverse intersections and
$D_q^{\mathrm{big}}$ the matrix of candidate bigons, put
\begin{equation}\label{eq:big-stat}
 h_q^{\mathrm{big}}:=\dim_{\F_2}C_q-2\rk_{\F_2}D_q^{\mathrm{big}};
\end{equation}
this is a number attached to a finite matrix, not the homology of a Floer complex.

At $q=5$ the curves meet in nine points (one interior, two clusters of four near $\gamma=\pi$), and the
test accepts exactly two candidate bigons with distinct vertices, so $h_5^{\mathrm{big}}=9-2\cdot2=5$;
this agrees with the nine intersections and two bigons in Smith's proof of \cite[Thm.~5.9]{Smith}.  The
generator counts depend on the parameters: at $q=11$ the matrix $D_{11}^{\mathrm{big}}$ has $19$ generators
and rank two at the parameters $(0.05,0.16,0.40)$, and $23$ generators and rank four at the parameters $(0.03,0.10,0.25)$, and at $q=13$ there are $27$ against $25$ generators; at $q=11$ the model has $240$ self-intersections, some less than $10^{-3}$ apart.  The value of $h_q^{\mathrm{big}}$
agrees at two independent parameter choices, at both of which the matrices $D_q^{\mathrm{big}}$ for
$q=5,7,11,13$ square to zero.  One non-generic earring parameter choice at $q=7$,
which collapses the $13$ generators to nine, was excluded.

\begin{computation}[Bigon ranks]\label{comp:law}
For the specified representatives and parameters,
\[
 h_5^{\mathrm{big}}=5,\qquad h_7^{\mathrm{big}}=7,\qquad
 h_{11}^{\mathrm{big}}=15,\qquad h_{13}^{\mathrm{big}}=17,
\]
and the matrices $D_q^{\mathrm{big}}$ square to zero over $\F_2$.  The value $5$ at $q=3$ in
Table~\ref{tab:family} is the computation of Hedden, Herald and Kirk \cite[\S11.6]{HHK2} from a
different decomposition; it is not an output of this family.
\end{computation}

\begin{table}[ht]
\centering
\small
\begin{tabular}{c|ccccc}
 $q$ & $\det K_q$ & $h_q^{\mathrm{big}}$ & $\dim_{\F_2}\Inat$ & difference & required rank change \\ \hline
 $3$ & $3$ & $5$ (external) & $5$ & $0$ & $0$ \\
 $5$ & $1$ & $5$ & $7$ & $-2$ & $-1$ \\
 $7$ & $1$ & $7$ & $9$ & $-2$ & $-1$ \\
 $11$ & $5$ & $15$ & $13$ & $+2$ & $+1$ \\
 $13$ & $7$ & $17$ & $15$ & $+2$ & $+1$ \\
\end{tabular}
\medskip
\caption{Bigon ranks and instanton dimensions.  The instanton column is~\cite[Thm.~1.1]{WuebbenII}.  The
column $h_q^{\mathrm{big}}$ for $q\ge5$ concerns the specified finite matrices; the $q=3$ entry is
external.  The last column is the change in matrix rank that a deformed differential on the same
generators would require.}
\label{tab:family}
\end{table}

A genuine deformation is a different object.  An immersed Floer theory has two ordered branch-jump
generators at each transverse self-intersection, of complementary degrees in a graded setup
\cite{AkahoJoyce}; a bounding cochain must satisfy the full Maurer--Cartan series, with repeated inputs
at every order; and square-zero of the deformed differential is necessary, though not sufficient, for
a proposed finite matrix.  The search below labels the underlying double point instead.  For a finite
set $S$ of self-intersections, a \emph{support}, let $D_q^{\mathrm{tab}}(S)$ add to
$D_q^{\mathrm{big}}$ the candidate triangle and quadrilateral terms through \emph{distinct} crossings
of $S$, and put
\begin{equation}\label{eq:tab-stat}
 h_q^{\mathrm{tab}}(S):=\dim_{\F_2}C_q-2\rk_{\F_2}D_q^{\mathrm{tab}}(S).
\end{equation}
Repeated insertions and higher operations are absent, and the accompanying Maurer--Cartan test uses
only the candidate monogon, self-bigon and distinct-input self-triangle terms.

\begin{computation}[Low-order candidate supports]\label{comp:cochains}
\begin{enumerate}
\item[(i)] At $q=5$, among the singleton and two-element subsets of $A_5$, the unique smallest
support with $h_5^{\mathrm{tab}}=7$ is $S=\{s_A,s_B\}$, with $(\gamma,\theta)\approx(0.03,1.27)$ and
$(3.06,4.98)$.  One candidate quadrilateral changes the $(x_4,y_6)$ bigon entry from $1$ to $0$; the
resulting matrix has the sole entry $(x_1,y_0)$ and squares to zero.
\item[(ii)] At $q=7$, among the tested singleton crossings passing the low-order test, one crossing near
$(0.042,5.405)$ has $h_7^{\mathrm{tab}}=9$.  Its candidate triangle changes one bigon entry from $1$ to
$0$; the resulting matrix has entries $(6,4)$ and $(7,5)$ and squares to zero.
\item[(iii)] At $q=11$ with the parameters $(0.05,0.16,0.40)$, $62$ singleton crossings pass the low-order
test and give matrix rank three.  Exactly $56$ of the resulting matrices square to zero; the square of each
of the other six has a single nonzero entry, namely $(9,15)$, $(9,16)$, $(0,14)$, and $(9,14)$ for the
remaining three.  At the parameters $(0.03,0.10,0.25)$ the same test returns $66$ singletons of the
target rank, of which $16$ fail square-zero and $50$ pass.
\end{enumerate}
\end{computation}

At $q=5$ the matrix $D_5^{\mathrm{big}}$ has rank two, and the pair $s_A,s_B$ of Figure~\ref{fig:pillow} is
the only tested support that lowers the rank to one, so $h_5^{\mathrm{tab}}=9-2=7$; minimality
holds only within this search.  At $q=7$ the rank three of $D_7^{\mathrm{big}}$ drops to two and
$h_7^{\mathrm{tab}}=13-4=9$; at the parameters $(0.07,0.16,0.40)$ the crossing is near
$(0.058,5.413)$ and gives the same entries.  This crossing is $S_{69}$, whose type-D deformation is the
four-arrow switch $b_{S_{69}}$ rather than a support labelled only by its double point, as in
Computation~\ref{comp:cochains}.    Resolving both local sectors at the selected $q=5$ and $q=7$ crossings and
recounting candidate bigons reproduces the matrices $D_q^{\mathrm{tab}}(S)$ for one sector choice; at $q=7$
Proposition~\ref{prop:q7-typeD} replaces this experiment by an exact calculation.
Computation~\ref{comp:cochains} proves neither existence nor uniqueness of a bounding cochain: the
candidates are not ordered branch jumps, the full Maurer--Cartan series is never tested, and neither $56$
nor $50$ counts bounding cochains.

\subsection{Methods}\label{ss:methods}
All statements of Sections~\ref{sec:smith} and~\ref{sec:khovanov} that rest on computation are finite
computations in exact arithmetic over $\F_2$, with the following exceptions, which are floating-point
computations: the geometry of the piecewise-linear curves of Section~\ref{sec:smith} (their self-intersections and
their crossings with the arc system, from which the encodings and the enumeration of smoothings are read off),
which uses numerical tolerances; the finite winding
test of \S\ref{sec:compute}, which evaluates winding functions on sampled regions; and the traced curves described
below.

\emph{Morphism homology.}  Morphism complexes in $\operatorname{Tw}(\mathcal B)$ are computed in
full, without truncating word length, by the mapping-cone reduction of
Lemma~\ref{lem:wrapped-reduction}, whose torsion condition is checked for each pair.  A second,
independent implementation computes $H\operatorname{Mor}(X,Y)$ from the structure of
$\operatorname{Mor}(X,Y)$ as a module over the polynomial ring $\F_2[H]$ of \cite{KWZ}; it
reproduces the dimensions of Propositions~\ref{prop:q7-typeD} and~\ref{prop:aux-closure}, and the
dimensions of Section~\ref{sec:khovanov} were computed with it.

\emph{Encodings.}  Curves given by their lifts to the planar cover are encoded as twisted complexes
by recording their crossings with the lifts of the arc system and the faces between consecutive
crossings, as in \S\ref{ss:explicit}.  The smoothed curves of Proposition~\ref{prop:q7-typeD} were encoded
with smoothing disks of radii $0.003$, $0.006$ and $0.009$, with the same result.  This encoder was compared with the invariants of $216$
rational tangles computed by Zibrowius's program \cite{Zibkht}: in every case
$\widetilde{\mathrm{BN}}(Q)$ is the straight lattice arc and $\widetilde{\Kh}(Q)$ the figure-eight
around it.  The identification of the earring curve $E_r$ with $\Phi\bigl(\widetilde{\Kh}(Q_{-r})\bigr)$ in
\S\ref{ss:kh-smith} was checked by comparing encodings at $49$ slopes, and the two sides
of~\eqref{eq:kwz-pairing} were compared at the $340$ slopes used for $q=5$ and the $342$ slopes used for $q=7$
in Corollary~\ref{cor:closures-full}; they agree at every slope.  The invariants $\widetilde{\mathrm{BN}}(T_q)$
and $\widetilde{\Kh}(T_q)$ for $q=3,5,\dots,21$ were computed with the same program, and the ranks
$\rk\widetilde{\Kh}(K_r(q);\F_2)$ with Khoca~1.5 \cite{Khoca}.

\emph{Traced curves.}  The perturbed curves of Smith for $q=5,7,11,13,17,19$ were traced numerically
from the equations of \cite[Prop.~4.1]{Smith}, composed with the lines of slopes $1/3$ and $1/q$
through $(0,0)$, for several perturbation parameters in $[0.005,0.02]$ and several step sizes, and
transported to negative perturbation parameter by the symmetry $\sm$ of Remark~\ref{rem:sign}.  For
$q\le13$ all traces at a given $q$ have strictly isomorphic encodings, isomorphic to $N_q$.  These are
floating-point computations and serve as corroboration only.

\emph{Searches.}  Strict isomorphism of loop-type complexes is decided by comparing the words of the
components up to rotation and reversal.  Edit distances are computed by an exact search over the ways
of cutting the word of one complex into intervals that reassemble to the word of the other; for
$q=17,19$ the search excludes edit distances at most $5$ and was stopped at $6$.  Chains of algebraic
switches are enumerated breadth first, up to strict isomorphism, from both ends, and the two sides are
compared.  Each search is exhaustive within the class it states, and no conclusion is drawn beyond
that class.

\subsection{Reproducibility}\label{ss:repro}
The programs that carry out the computations of this paper are included as ancillary files with its
arXiv submission, and they are also available in the public repository
\begin{center}
\url{https://github.com/bwuebben/pillowcase-bounding-cochain}\,.
\end{center}
In the ancillary files the programs for Section~\ref{sec:smith} and \S\ref{sec:compute} are in the directory
\path{pillowcase/}, and those for Section~\ref{sec:khovanov} and the rest of this appendix are in \path{code/}; in the repository these are the
directories \path{pillowcase/} and \path{paper2b/code/}. Each program prints the numbers of the statements it checks next to the computed values and reports whether the two agree, and the file \path{README.md} in each directory
lists the statements that each program checks; there, oriented and unoriented smoothings are called census
and opposite smoothings, and connector--inherited self-intersections connector--main ones. The programs use only
the Python standard library. The invariants
computed with Zibrowius's program~\cite{Zibkht} and the ranks computed with Khoca~\cite{Khoca} are included as
data, together with their input files, so the checks run without these programs. The instanton dimensions used
throughout are~\cite[Thm.~1.1]{WuebbenII} and Proposition~\ref{prop:aux-closure}, and they are independent of all
of this code.
